\subsection{On the Geometry of the String Landscape and the Swampland --- arXiv:hep-th/0605264}

\textbf{Authors:} Hirosi Ooguri, Cumrun Vafa

\paragraph{主な主張}
量子重力のモジュライ空間に関する幾何学的予想を提案し、無限距離極限では無限個の軽い状態の塔が現れるという条件を弦理論の例から動機付ける\cite{Ooguri:2006in}。

\paragraph{新規性}
モジュライ空間の距離・曲率・位相の性質を、低エネルギー理論が量子重力と両立するための条件として結び付ける。

\paragraph{位置付け}
距離予想の原点となる文献であり、大きなモジュラス変位における有効理論の適用範囲を考える基礎。

\paragraph{コメント（読書メモ）}
Distance conjectureに関する原論文。
